% Miguel - Ejercicio 3, incisos a-c: traducción e inserción de tuplas.

\section{Ejercicio 3. Modelo relacional e inserción de tuplas}

En el modelo E/R del enunciado, \emph{A} tiene la llave compuesta
$(a_1,a_2)$: el atributo subrayado $a$ se descompone en $a_1$ y $a_2$.
La llave de \emph{B} es $b$, y $ab_1$ pertenece a la relación \emph{ab}.
Se considera participación parcial de ambos lados en las cuatro
cardinalidades, como pide el inciso~a.

\subsection{Inciso a. Traducción para cada cardinalidad}

La Tabla~\ref{tab:ej3-traduccion} presenta las relaciones resultantes.
El subrayado continuo identifica la llave primaria completa; las flechas
indican las llaves foráneas y su relación de destino. Dentro de cada esquema
se colocan primero las llaves primarias, después los atributos propios y al
final las llaves foráneas, conservando de izquierda a derecha el orden de los
atributos del modelo E/R.

\begin{table}[htbp]
  \centering
  \caption{Traducción del modelo E/R según la cardinalidad de \emph{ab}.}
  \label{tab:ej3-traduccion}
  \small
  \begin{tabularx}{\textwidth}{@{}l >{\raggedright\arraybackslash}X@{}}
    \toprule
    \textbf{Cardinalidad} & \textbf{Modelo relacional e integridad referencial} \\
    \midrule
    $M:N$ &
      \begin{tabular}[t]{@{}l@{}}
        $\mathrm{A}(\underline{a_1,a_2},a_3)$;
        $\mathrm{B}(\underline{b},b_1)$ \\
        $\mathrm{AB}(\underline{a_1,a_2,b},ab_1)$ \\
        $\mathrm{AB}(a_1,a_2)\to\mathrm{A}(a_1,a_2)$;
        $\mathrm{AB}.b\to\mathrm{B}.b$
      \end{tabular} \\
    $1:N$ &
      \begin{tabular}[t]{@{}l@{}}
        $\mathrm{A}(\underline{a_1,a_2},a_3)$;
        $\mathrm{B}(\underline{b},b_1,ab_1,a_1,a_2)$ \\
        $\mathrm{B}(a_1,a_2)\to\mathrm{A}(a_1,a_2)$
      \end{tabular} \\
    $N:1$ &
      \begin{tabular}[t]{@{}l@{}}
        $\mathrm{A}(\underline{a_1,a_2},a_3,ab_1,b)$;
        $\mathrm{B}(\underline{b},b_1)$ \\
        $\mathrm{A}.b\to\mathrm{B}.b$
      \end{tabular} \\
    $1:1$ &
      \begin{tabular}[t]{@{}l@{}}
        $\mathrm{A}(\underline{a_1,a_2},a_3)$;
        $\mathrm{B}(\underline{b},b_1)$ \\
        $\mathrm{AB}(\underline{a_1,a_2},b,ab_1)$;
        $\operatorname{UNIQUE}(\mathrm{AB}.b)$ \\
        $\mathrm{AB}(a_1,a_2)\to\mathrm{A}(a_1,a_2)$;
        $\mathrm{AB}.b\to\mathrm{B}.b$
      \end{tabular} \\
    \bottomrule
  \end{tabularx}
\end{table}

En $M:N$ una fila de \emph{AB} identifica una pareja de entidades; su llave
primaria $(a_1,a_2,b)$ impide registrar dos veces esa pareja. En $1:N$ se
colocan $(a_1,a_2)$ y $ab_1$ en \emph{B}, que es el lado $N$; en $N:1$ se
colocan $b$ y $ab_1$ en \emph{A}, que ahora es el lado $N$. Como la
participación es parcial, la llave foránea de una entidad sin pareja puede
quedar sin valor. En el caso de la llave compuesta, sus dos componentes deben
estar ausentes juntos o formar una pareja existente en \emph{A}; $ab_1$ queda
sin valor si no existe una asociación.

Para $1:1$ se conserva \emph{AB} como relación independiente. La llave
primaria $(a_1,a_2)$ impide que una fila de \emph{A} tenga dos parejas y la
restricción $\operatorname{UNIQUE}(b)$ impide que una fila de \emph{B} tenga
dos parejas. La ausencia de una fila en \emph{AB} representa la participación
parcial de cualquiera de las dos entidades. Esta traducción también deja
disponible la relación \emph{AB} que utiliza el inciso~d.

\subsection{Inciso b. Inserciones con cardinalidad $M:N$}

Se parte de $\mathrm{A}(a_1,a_2,a_3)$ con las llaves
$(2,\text{'ww'})$, $(4,\text{'xx'})$, $(6,\text{'yy'})$ y
$(8,\text{'zz'})$, y de $\mathrm{B}$ con las llaves
$17,27,37,47,57$. Cada tupla propuesta para
$\mathrm{AB}(a_1,a_2,b,ab_1)$ debe tener una pareja $(a_1,a_2)$ presente
en \emph{A}, un valor $b$ presente en \emph{B}, un entero en $ab_1$ y una
llave primaria $(a_1,a_2,b)$ no repetida.

\begin{table}[htbp]
  \centering
  \caption{Evaluación de los conjuntos del inciso 3.b, en el orden propuesto.}
  \label{tab:ej3-b}
  \small
  \begin{tabularx}{\textwidth}{@{}l l >{\raggedright\arraybackslash}X@{}}
    \toprule
    \textbf{Conjunto} & \textbf{Completo} & \textbf{Justificación} \\
    \midrule
    i & No & La quinta tupla repite la tercera,
      $(4,\text{'xx'},27,15)$; se repite la llave primaria
      $(4,\text{'xx'},27)$. \\
    ii & No & Desde la primera tupla, $(17,\text{'zz'},2,\text{'m'})$,
      $(17,\text{'zz'})$ no es llave de \emph{A}, $2$ no es llave de
      \emph{B} y $ab_1=\text{'m'}$ no es entero. \\
    iii & No & La primera tupla usa $(2,\text{'a'})$ como llave foránea;
      \emph{A} contiene $(2,\text{'ww'})$. De igual modo, las cadenas
      \texttt{'b'}, \texttt{'c'} y \texttt{'d'} son valores de $a_3$,
      no componentes $a_2$ de las llaves de \emph{A}. \\
    iv & Sí & Las cinco parejas $(a_1,a_2)$ y los cinco valores $b$
      existen, $ab_1$ es entero y las cinco llaves $(a_1,a_2,b)$
      son distintas. \\
    \bottomrule
  \end{tabularx}
\end{table}

Por tanto, \textbf{únicamente el conjunto iv se inserta por completo}.
Para cada conjunto rechazado se propone a continuación una alternativa de
cinco tuplas válidas, en el orden
$\mathrm{AB}(a_1,a_2,b,ab_1)$:

\begin{enumerate}[label=\roman*.,leftmargin=*,itemsep=0.2em]
  \item \texttt{(8,'zz',17,5); (6,'yy',57,10);}\\
    \texttt{(4,'xx',27,15); (2,'ww',37,20);}\\
    \texttt{(4,'xx',47,15).}
    Se cambia la última llave $b=27$ por $b=47$.
  \item \texttt{(8,'zz',17,2); (6,'yy',27,4);}\\
    \texttt{(4,'xx',37,6); (2,'ww',47,8);}\\
    \texttt{(4,'xx',57,10).}
    Se usan llaves de \emph{A} y \emph{B} existentes y enteros en $ab_1$.
  \item \texttt{(2,'ww',17,23); (4,'xx',27,24);}\\
    \texttt{(6,'yy',37,25); (8,'zz',47,26);}\\
    \texttt{(2,'ww',57,27).}
    Se sustituye cada valor de $a_3$ por el $a_2$ correspondiente.
\end{enumerate}

\subsection{Inciso c. Inserciones con cardinalidad $1:N$}

En esta traducción las tuplas se insertan en
$\mathrm{B}(b,b_1,ab_1,a_1,a_2)$. Se conservan las cuatro tuplas de
\emph{A} indicadas en el inciso~b y se evalúa cada conjunto desde una
\emph{B} vacía: las cinco tuplas de \emph{B} citadas allí describían el
estado del modelo $M:N$ y no son inserciones previas de este inciso.
Se exige que $b$ sea único, que los tipos sigan el orden
entero--cadena--entero--entero--cadena y que la pareja $(a_1,a_2)$
referencie a \emph{A}. Varios valores de $b$ pueden referenciar a la misma
fila de \emph{A}, pues \emph{B} es el lado $N$.

\begin{table}[htbp]
  \centering
  \caption{Evaluación de los conjuntos del inciso 3.c, en el orden propuesto.}
  \label{tab:ej3-c}
  \small
  \begin{tabularx}{\textwidth}{@{}l l >{\raggedright\arraybackslash}X@{}}
    \toprule
    \textbf{Conjunto} & \textbf{Completo} & \textbf{Justificación} \\
    \midrule
    i & No & En $(2,\text{'f'},57,\text{'zz'},10)$ la llave foránea
      tendría $a_1=\text{'zz'}$ y $a_2=10$, de tipos incorrectos.
      Además, la quinta tupla repite $b=2$ de la primera. \\
    ii & No & En $(57,8,\text{'zz'},\text{'f'},2)$, $b_1$ y $a_2$
      deberían ser cadenas, mientras $ab_1$ y $a_1$ deberían ser enteros.
      El mismo cambio de posiciones afecta a las demás tuplas. \\
    iii & No & En $(17,\text{'f'},8,\text{'zz'},1)$ los componentes
      de la llave foránea tienen los tipos invertidos. También se repiten
      las llaves primarias $b=27$ y $b=17$ dentro del conjunto. \\
    iv & Sí & Los valores de $b$ son distintos, $ab_1$ es entero y las
      llaves foráneas son $(2,\text{'ww'})$, $(4,\text{'xx'})$,
      $(6,\text{'yy'})$, $(8,\text{'zz'})$ y $(4,\text{'xx'})$;
      todas existen en \emph{A}. \\
    \bottomrule
  \end{tabularx}
\end{table}

Así, \textbf{únicamente el conjunto iv se inserta por completo}.
Los siguientes conjuntos de cinco tuplas corrigen, respectivamente, los
casos i, ii y iii, siempre en el orden
$\mathrm{B}(b,b_1,ab_1,a_1,a_2)$:

\begin{enumerate}[label=\roman*.,leftmargin=*,itemsep=0.2em]
  \item \texttt{(57,'f',10,8,'zz'); (47,'g',15,6,'yy');}\\
    \texttt{(37,'h',10,4,'xx'); (27,'i',15,2,'ww');}\\
    \texttt{(17,'j',20,6,'yy').}
  \item \texttt{(57,'f',2,8,'zz'); (47,'g',4,6,'yy');}\\
    \texttt{(37,'h',6,4,'xx'); (27,'i',8,2,'ww');}\\
    \texttt{(17,'j',6,6,'yy').}
  \item \texttt{(17,'f',1,8,'zz'); (27,'g',2,6,'yy');}\\
    \texttt{(37,'h',3,4,'xx'); (47,'i',4,2,'ww');}\\
    \texttt{(57,'j',5,6,'yy').}
\end{enumerate}

En las tres alternativas los cinco valores de $b$ son distintos, los
atributos respetan sus tipos y cada pareja $(a_1,a_2)$ existe en
\emph{A}; por ello se pueden insertar las cinco tuplas de cada conjunto.
